\documentclass[12pt, letterpaper]{article} 

\usepackage{amsmath, amssymb, amsthm, enumerate, graphicx}

\title{MATH 114 HW 3}
\author{Dani Nguyen}
\date{Winter 2025}

\begin{document}
\maketitle
\graphicspath{ {./} }

\noindent \textbf{Question 1} \\
Let $k = 1$ be the class of companies that issue dividends and $k=2$ be the class 
of companies that didn't. \\
Let $X_1$ be the percent profit of $k=1$ and $X_2$ be the percent profit of $k=2$ \\
We have
\begin{align*}
    X_1 \sim N(10,36) && X_2 \sim N(0,36)
\end{align*}
\begin{align*}
    P(k=1) = .8
\end{align*}
\begin{align*}
    P(k=1 | X) &= \frac{P( y = 1 ) P(X | y = 1)}{P( y = 1 ) P(X | y = 1) + P( y = 2 ) P(X | y = 2)} \\
    &= \frac{.8 f_1(X)}{.8 f_1(X) + .2 f_2(X)} \\
    &= \frac{.8 \frac{1}{\sqrt{2\pi\sigma^2}} \exp (-\frac{1}{2\sigma^2}(X-\mu_1)^2)}{.8 \frac{1}{\sqrt{2\pi\sigma^2}} \exp (-\frac{1}{2\sigma^2}(X-\mu_1)^2) + .2 \frac{1}{\sqrt{2\pi\sigma^2}} \exp (-\frac{1}{2\sigma^2}(X-\mu_2)^2)} \\
    &= \frac{.8 \exp (-\frac{1}{2\cdot 36}(4-10)^2)}{.8 \exp (-\frac{1}{2\cdot 36}(4-10)^2) + .2 \exp (-\frac{1}{2\cdot 36}(4-0)^2)} \\
    &= \frac{.8 \exp (-\frac{1}{72}(-6)^2)}{.8 \exp (-\frac{1}{72}(-6)^2) + .2 \exp (-\frac{1}{72}(4)^2)} \\
    &= \frac{.8 \exp (-1/2)}{.8 \exp (-1/2) + .2 \exp (-2/9)} \\
    &= 0.7519
\end{align*}
\noindent \textbf{Question 2} \\
When K=1, KNN fit perfectly interpolates the training observation which gives 
zero error on the training data set. This means the test error rate is 36\%. 
Therefore, logistic regression is better. \\

\noindent \textbf{Question 3} 
\begin{enumerate}[a.]
    \item We split our data according to $k$ subsets, and then we save each 
    $k$'th set for testing, and training on the remaining $k-1$ sets combined.
    \item A drawback of the validation set approach is that it is very sensitive 
    to which data is included in the training, it is however less computationally 
    expensive because we only have to fit one model. For the LOOCV, we set $k=n$
    and use most of the data for training, saving one observation for validation.
    This is very computationally expensive compared to validation sets, but has
    far less bias.
\end{enumerate}

\noindent \textbf{Question 4} 

\noindent \textbf{Question 5} 
\begin{enumerate}
    \item 
\end{enumerate}

\end{document}